% ==================================================================
% APPENDIX A6 — Methodology + benchmark achievability + age-viability
%
% FINAL DRAFT on freeze-checkpoint numbers (2026-04-19).
% This appendix is referenced as \ref{app:eps_methodology}
% throughout A1-A5 and §15.5.
%
% Structure:
%   A6.1  Common background module: scope and role
%   A6.2  What was and was not recomputed under the module
%   A6.3  Freeze-checkpoint joint band: provisional but stable
%   A6.4  Age-viability ambiguity under Hubble mapping (viability
%         note, NOT a retained probe)
%   A6.5  Benchmark achievability for the retained band (only)
%   A6.6  Excluded probes: audit status and non-silent-dismissal
%   A6.7  Methodological limitations and path forward
%
% Character (GPT r17-r21):
%   A6 is the binding glue. Without it, A1-A5 are disconnected
%   per-channel extractions. A6 makes the retained-five-probe
%   analysis a coherent package.
%
% Target audience:
%   - readers who need to understand why the retained band is
%     "provisional but stable"
%   - reviewers who will ask about age tension
%   - anyone auditing excluded probes or asking "why aren't BAO
%     and S_8 headline?"
%
% NO integration into preprint yet.
% ==================================================================

\section{Methodology of the common-background retained-pipeline,
benchmark achievability, and the age-viability ambiguity}
\label{app:eps_methodology}

This appendix collects the methodological scaffolding for the
retained-five-probe analysis of \S15.5. Its purpose is threefold:
(i) to specify the common background module that A1--A5 all
invoke, (ii) to state the current status of the retained joint
band as ``provisional but stable under the first common-background
recalc,'' and (iii) to discuss two issues that belong at this
cross-channel level rather than in any single per-probe appendix
---namely, the age-viability ambiguity under the Hubble mapping
and the benchmark achievability for the retained band. The
appendix closes with an audit of the excluded probes and a
discussion of the methodological limitations still open after the
freeze checkpoint.

\subsection*{A6.1\quad The common background module: scope and
role}

All five retained per-probe extraction appendices (A1--A5) invoke
a single background module, \texttt{ect\_background.py}, which
functions as the single source of truth for the
$\varepsilon$-deformed background quantities used across channels.
Specifically, the module provides:
\begin{itemize}
\item the dimensionless Hubble rate
$E_{\rm ECT}(z,\,\varepsilon)$ from
\eqref{eq:eps_friedmann_A1};
\item the Hubble rate
$H(z,\,\varepsilon) = H_0 \cdot E_{\rm ECT}(z,\,\varepsilon)$
at a user-supplied $H_0$;
\item the cosmic age
$t_0(\varepsilon)
  = \int_0^{\infty}\frac{dz}{(1+z)\,H(z,\,\varepsilon)}$
evaluated with a multi-segment quadrature to avoid
round-off in the deep-matter-era integration range;
\item the lookback time and age at arbitrary redshift,
$t(z,\,\varepsilon)$, used in A2 to construct the
stellar-assembly time-budget factor;
\item the standard cosmological parameter constants
($\Omega_m$, $\Omega_\Lambda$, $\Omega_r$, $\Omega_b$,
$z_*$, $z_{\rm drag}$, $c$, $H_0^P$) adopted uniformly
across channels.
\end{itemize}
At $\varepsilon = 0$ the module reproduces the published
$\Lambda$CDM benchmarks ($t_0 \approx 13.79~{\rm Gyr}$) at
better than $1\%$ precision once the integration range fix is
applied. The module does not implement recombination-era or
Boltzmann-level physics, nor does it implement the closure-level
ECT background; it is a proxy that allows all per-channel
extractions to share the same low- and late-$z$ numerical
scaffolding.

\subsection*{A6.2\quad What was and was not recomputed under the
common module}

In preparing the present checkpoint, each retained channel was
re-expressed against the common module of A6.1. The level of
physical completion, however, differs by channel. The following
summary is the normative table: it is the reference that the
``provisional but stable'' language of
\S\ref{subsec:eps_retained_rebuilt} and
\S\ref{subsec:eps_joint_band_rebuilt} points back to.
\begin{center}
\begin{tabular}{|l|l|l|}
\hline
Channel & Status after freeze &
Comment \\
\hline
Hubble + $r_s$ (A1) &
stable &
No change from pre-recalc. \\
& & Weak $t_0$ coupling at leading order. \\
\hline
JWST (A2) &
materially revised; &
Time-budget factor added; central $\varepsilon$ \\
& appendix-level defence required &
 shifted from $0.0425$ to $0.0449$. \\
\hline
Cosmic Chronometers (A3) &
stable &
Imported background; $H(z)$ shape probe. \\
& & One-sided upper under $\varepsilon \geq 0$. \\
\hline
$f\sigma_8$ RSD (A4) &
stable after grid-fix &
Brent-refined $1\sigma$ upper
$0.0396$. \\
& & Raw negative best-fit explicit in A4.5. \\
\hline
ISW (A5) &
provisional, proxy-level &
$\kappa_{\rm ISW}$ \emph{not} recalibrated \\
& & on common background. \\
\hline
\end{tabular}
\end{center}
\paragraph{Net effect.}
Two channels carried no numerical change (Hubble, CC); two
channels recovered or stabilised at their previously published
values after re-grounding on the common module (fσ_8, ISW); one
channel (JWST) received a materially new physical correction
(time-budget factor). The net impact on the retained joint band
is discussed in A6.3.

\subsection*{A6.3\quad The freeze-checkpoint joint band: provisional
but stable}

The retained joint $1\sigma$ band after the common-background
recalc is
\begin{equation}
\varepsilon
\;\in\;
[\,0.0296,\ 0.0376\,]
\quad (1\sigma)
\label{eq:joint_band_A6}
\end{equation}
with the $2\sigma$ extension $[\,0.0207,\,0.0425\,]$. The binding
edges are the JWST lower-$1\sigma$ (A2, now $0.0296$ after the
time-budget correction) and the Hubble + $r_s$ upper-$1\sigma$
(A1, stable at $0.0376$). The $f\sigma_8$ upper-$1\sigma$
($0.0396$, Brent-refined) is $0.002$ looser than the Hubble edge
and does not bind; the cosmic-chronometer ($0.087$) and ISW
($0.043$) upper edges are broader still.

\paragraph{Interpretation.}
This band is numerically stable enough to support appendix
drafting, but not yet final enough for main-text updating. The
overall numerical match of \eqref{eq:joint_band_A6} with the
previous pre-recalc band $[\,0.029,\,0.036\,]$ (to three
significant figures) indicates that the first common-background
recalc preserves the retained-five-probe intersection at its
previous location, within the precision of the present
pipelines. This suggests numerical stability of the intersection
under the first common-background recalc, but does not by itself
establish the final retained result. Final confirmation is
pending completion of two pieces of work: the fully rewritten
JWST time-budget defence at appendix level (A2) and the
ECT-native kernel recalibration for ISW (A5). Both are listed as
methodology limitations in A6.7.

\paragraph{What the band does and does not claim.}
The band is the intersection of five independent $1\sigma$
intervals on the uniform-$\varepsilon$ diagnostic layer. It does
not claim to be a ``measurement of $\varepsilon$'': the
uniform-$\varepsilon$ ansatz is itself a projection of a more
general closure-level $\varepsilon(z)$ history onto a single
effective parameter (cf.\ A1.7(iii)). The band is also not a
direct fit to a first-principles ECT closure; it is the
compatibility region under which the effective parameter
simultaneously accommodates five independent observational
channels at the present proxy level.

\subsection*{A6.4\quad Age-viability ambiguity under Hubble mapping
(viability note)}

The $\varepsilon$-deformed background of A6.1 computes the cosmic
age $t_0(\varepsilon)$ self-consistently. At the retained-band
central value $\varepsilon = 0.032$, $t_0$ is approximately
$13.46~{\rm Gyr}$ when the age integral is evaluated at the
baseline Hubble constant $H_0 = H_0^P = 67.4~{\rm km\,s^{-1}\,Mpc^{-1}}$,
and approximately $12.52~{\rm Gyr}$ when the same age integral is
evaluated at the \emph{inferred} Hubble constant
$H_0^{\rm inferred}(\varepsilon = 0.032) \approx 72.55~{\rm km\,s^{-1}\,Mpc^{-1}}$
from the Hubble + $r_s$ mapping of \eqref{eq:H0_inferred_eps}.

\paragraph{Comparison with observational age anchor.}
The Valcin et al.\ 2021 globular-cluster age analysis
\cite{Valcin2021} gives a largely model-independent lower bound
$t_0 \geq 13.50 \pm 0.27~{\rm Gyr}$ on the cosmic age. The two
mappings above yield qualitatively different comparisons:
\begin{itemize}
\item baseline-background mapping: $13.46~{\rm Gyr}$ is
   within $0.15\sigma$ of the anchor --- acceptable;
\item inferred-$H_0$ mapping: $12.52~{\rm Gyr}$ is
   approximately $3.6\sigma$ below the anchor --- problematic.
\end{itemize}
\paragraph{Interpretation.}
A nontrivial age-viability ambiguity appears once the Hubble
extraction is combined with the uniform-$\varepsilon$ background
interpretation. Whether this should be read as a genuine
consistency tension or as an artefact of how the inferred $H_0$
is mapped back into the background solution requires dedicated
treatment that is \emph{not} performed at the uniform-$\varepsilon$
diagnostic layer. Specifically, three interpretations are
currently open:
\begin{itemize}
\item the baseline-background mapping is the physically correct
  reading of ``cosmic age,'' because the ECT closure only
  shifts the \emph{apparent} local Hubble constant without
  altering the absolute cosmic expansion history;
\item the inferred-$H_0$ mapping is the physically correct
  reading, in which case the uniform-$\varepsilon$ ansatz
  becomes inadequate near the retained-band central value and
  a closure-level $\varepsilon(z)$ must decrease toward the
  present;
\item the globular-cluster age anchor is subject to calibration
  systematics at the $\sim 1~{\rm Gyr}$ level, in which case
  the observed tension is consistent with measurement
  uncertainty.
\end{itemize}
None of these three interpretations is fixed by the
Hubble + $r_s$ channel alone. The present appendix does not
resolve the ambiguity.

\paragraph{Status.}
Age is \emph{not} a retained probe. It is \emph{not} a new formal
open problem (OP-age-consistency is not introduced). It is
recorded here as a viability/self-consistency note. Its
resolution is part of the eventual closure-level ECT derivation
of $\varepsilon(z)$ and $H_0$ from first principles (cf.\
\S\ref{subsec:eps_outlook_rebuilt}), not a constraint to be added
at the uniform-$\varepsilon$ level.

\subsection*{A6.5\quad Benchmark achievability for the retained
band (only)}

``Benchmark achievability'' refers to the question: within the
parameter space of the ECT ordered-branch closure, do there
exist configurations $(\beta,\,\mu,\,\kappa,\,\phi_0,\,\ldots)$
that yield a uniform-$\varepsilon$ projection in the retained
band \eqref{eq:joint_band_A6}? The question is posed purely at
the effective-diagnostic-layer level: can the retained band be
reached?

\paragraph{Scope.}
The target of the benchmark achievability analysis is the
retained band $[\,0.0296,\,0.0376\,]$ (\S A6.3) and nothing else.
In particular, the age-viability ambiguity of A6.4 is not mixed
into the benchmark target --- it remains a separate viability
consideration at the closure-derivation level, not an additional
constraint at the diagnostic level. This separation is
deliberate: conflating the two would confuse the purpose of the
benchmark achievability analysis and would cause the target
interval to depend on which interpretation of the age-viability
ambiguity one adopts.

\paragraph{What is required.}
A benchmark achievability demonstration must show that, at the
closure level of the condensate theory
(three-stage condensate closure at
$u_0 \sim \phi_0 \sim \bar{M}_{\rm Pl}$,
$v_2 \sim 246~{\rm GeV}$, and $v_{\rm gal} \sim$ kpc; cf.\
\S\ref{subsec:eps_outlook_rebuilt}), at least one parameter
configuration yields an effective $\varepsilon$ at late time in
the range $[\,0.0296,\,0.0376\,]$ when projected onto the
uniform-$\varepsilon$ diagnostic. No claim is made here that
\emph{every} closure-level history produces a uniform-$\varepsilon$
projection in this range, only that the range is reachable.

\paragraph{Current status.}
Benchmark achievability for the retained band is an open target:
it has not been demonstrated in the present preprint. The
closure-level derivation of $\varepsilon(z)$ from three-stage
condensate dynamics is listed as open problem OP-Hubble-derive
in \S\ref{subsec:eps_outlook_rebuilt}, and an explicit check of
whether its output projection onto uniform-$\varepsilon$ lies in
the retained band is part of that open program.

\subsection*{A6.6\quad Excluded probes: audit status and
non-silent-dismissal}

Several probes that would naively contribute to a combined
multi-messenger analysis are explicitly excluded from the
retained joint band of \S\ref{subsec:eps_retained_rebuilt}. It
is important that their exclusion not be read as a silent
dismissal: the reasons for exclusion differ by probe, and each
remains an open item at the closure-level program.
\begin{center}
\begin{tabular}{|l|l|l|}
\hline
Probe & Status & Nature of exclusion \\
\hline
BAO (DESI 2024) & excluded; audited & Raw best-fit
$\varepsilon \approx -0.014$ \\
& & under simplified pipeline; \\
& & diagonal errors, shape-only. \\
\hline
$A_{\rm lens}$ CMB lensing & excluded; proxy only & Linearised
without Boltzmann; \\
& & theoretical systematic not included. \\
\hline
$S_8$ KiDS-1000 & excluded; outside sector & Raw $1\sigma$
entirely negative; \\
& & measures $S_8$ tension, not ECT $\varepsilon$. \\
\hline
$\sigma_8$ clusters & excluded; outside sector & Same as $S_8$. \\
\hline
SN Ia & excluded; not reinstated & Was a mock channel; permanently
removed. \\
\hline
Age $t_0$ & excluded; viability only & See A6.4. \\
\hline
\end{tabular}
\end{center}

\paragraph{Key distinction: excluded $\neq$ silently dismissed.}
All probes above have been audited at the level of the common
background module (A6.1). Their exclusion from the headline joint
band reflects a methodological judgement: at the present proxy
level, either (a) the probe's central value and/or $1\sigma$
lies in the unphysical region $\varepsilon < 0$, (b) the probe
measures a quantity whose primary sensitivity lies outside the
uniform-$\varepsilon$ sector, or (c) the probe's pipeline is
proxy-level enough that including it in the headline intersection
would be misleading. In no case has a probe been dropped simply
because it disagrees with the retained band. Any probe could in
principle re-enter the retained analysis under a closure-level
full pipeline; but none enters the present freeze-checkpoint
headline.

\subsection*{A6.7\quad Methodological limitations and path forward}

The following are known limitations of the present
freeze-checkpoint retained pipeline, to be addressed in future
work. They are stated here so that future revisions can be
tracked against them:
\begin{itemize}
\item the common background module is a $\Lambda$CDM-proxy,
  not a closure-level ECT background;
\item the ISW kernel $\kappa_{\rm ISW}$ has not been recalibrated
  on the common background (A5.3, A5.7(ii));
\item the JWST time-budget factor has been defended at robustness
  level for two axes (A2.4) but is still a proxy for the full
  halo-to-stellar-mass pipeline;
\item off-diagonal covariance has not been propagated for the
  $f\sigma_8$ RSD compilation (A4.7(ii)) and the cosmic
  chronometers compilation (A3.7(i));
\item the joint band \eqref{eq:joint_band_A6} is not yet
  promoted from provisional to final; main-text updating is
  deferred until appendix-package completion and review;
\item the age-viability ambiguity of A6.4 is not resolved at
  the present layer.
\end{itemize}

\paragraph{Path forward.}
The immediate step beyond the present freeze checkpoint is the
closure-level derivation of $\varepsilon(z)$ from the three-stage
condensate program, which is the content of open problem
OP-Hubble-derive (\S\ref{subsec:eps_outlook_rebuilt}). The
per-channel appendices A1--A5 and the present methodology
appendix A6 together define the scaffolding against which the
closure-level output will eventually be benchmarked.

% ==================================================================
% END Appendix A6 (methodology + benchmark + age-viability).
% ==================================================================
